% -----------------------------------------------------------------------------
%  entries/q025.tex
%  Issued 21/10/2026. First entry of the GKP / magic / hypergraph sequence
%  (Q025 GKP code; Q026 hyperedge as logical CCZ; Q027 magic-state injection with
%  hypergraph states; Q028 Wigner negativity versus logical magic).
%  All items solved before issue. Checked numerically: logical error rates of (c)
%  for d = 2..5 (exact series vs leading term); the logical phase gates of (d)
%  for d = 2..5.
%  22/10/2026: (c) reframed without numerical values (closed-form d-scaling only);
%  the Clifford argument moved from (d)(iii) to a new item (e), which now states
%  the definition of a Gaussian unitary it needs.
% -----------------------------------------------------------------------------
\begin{entry}{Q025}{A qudit inside an oscillator: the GKP code}

\entrymeta{
  topic   = {Bosonic codes; GKP code; qudit Pauli and Clifford groups; Gaussian operations},
  est     = {$\sim$2--3 hours including the write-up},
  origin  = Original,
  source  = {Original problem on the code of D.~Gottesman, A.~Kitaev and J.~Preskill,
    \emph{Encoding a qubit in an oscillator}, Phys.~Rev.~A \textbf{64}, 012310 (2001)},
  url     = https://arxiv.org/abs/quant-ph/0008040,
  prereq  = {Displacement operators and BCH with a central commutator (toolbox T7);
    the additive Gaussian noise of \qref{Q023}},
  status  = drafted,
}

\begin{problem}
\emph{Context.} An oscillator has an infinite-dimensional Hilbert space, and the noise it
suffers in optics and microwave cavities (loss, heating, imperfect Gaussian gates) acts, to a
good approximation, as small random shifts of $q$ and $p$. The code of Gottesman, Kitaev and
Preskill (GKP) hides a $d$-level system inside the oscillator so that every shift smaller
than a fixed size can be detected and undone. It is the error-correcting layer of current
photonic architectures. This entry builds the code from its stabilizers, finds which shifts
it corrects and at what noise level, and shows that its logical Clifford gates are all
Gaussian. Where the remaining, non-Clifford gates come from is the subject of the next
entries.

\emph{Setting.} Take one mode with $[\hat q,\hat p]=i$, $\hbar=1$, and position
eigenstates $\ket{q}$. For real $a,b$ define the shifts
\begin{equation}
  X(a):=e^{-ia\hat p},\qquad Z(b):=e^{ib\hat q},
  \qquad X(a)\ket{q}=\ket{q+a},\qquad Z(b)\ket{q}=e^{ibq}\ket q ,
  \label{q025:shifts}
\end{equation}
so that $Z(b)$ shifts $\hat p$ by $b$. By BCH with a central commutator \usetool{T7},
\begin{equation}
  Z(b)\,X(a)=e^{iab}\,X(a)\,Z(b).
  \label{q025:weyl}
\end{equation}
Fix an integer $d\ge2$, let $\ell:=\sqrt{2\pi/d}$ and $\omega:=e^{2\pi i/d}$, and define
the \emph{stabilizers} and \emph{logical Pauli operators}
\begin{equation}
  S_q:=X(d\ell),\quad S_p:=Z(d\ell),\qquad \bar X:=X(\ell),\quad \bar Z:=Z(\ell).
  \label{q025:code}
\end{equation}
The \emph{code space} consists of the (unnormalisable) states fixed by $S_q$ and $S_p$. For
$d=2$ this is the square GKP qubit, with $\ell=\sqrt\pi$. Treat the states formally, as in
\qref{Q022}.

\begin{enumerate}[label=\textbf{(\alph*)},leftmargin=*]
  \item \emph{A qudit inside the oscillator.} Show that $S_q$ and $S_p$ commute with each
        other and with $\bar X$ and $\bar Z$, that $\bar X^{d}=S_q$, $\bar Z^{d}=S_p$, and
        that $\bar Z\bar X=\omega\,\bar X\bar Z$. Show that the states
        \begin{equation}
          \ket{\bar\jmath}:=\sum_{k\in\ZZ}\ket{q=(dk+j)\ell},\qquad j\in\ZZ_d ,
          \label{q025:codewords}
        \end{equation}
        lie in the code space and satisfy $\bar Z\ket{\bar\jmath}=\omega^{j}\ket{\bar\jmath}$
        and $\bar X\ket{\bar\jmath}=\ket{\overline{j+1}}$.

  \item \emph{Which shifts are correctable.} Let a shift error $E=X(u)Z(v)$, with $u,v$
        real, act on a code state $\ket{\bar\psi}$.
        \begin{enumerate}[label=(\roman*)]
          \item Show that the eigenvalues of $S_q$ and $S_p$ on $E\ket{\bar\psi}$
                determine $v$ and $u$ modulo $\ell$, and nothing more.
          \item The correction applies $R:=X(-\bar u)\,Z(-\bar v)$, where $\bar u$ and
                $\bar v$ are the representatives of $u$ and $v$ modulo $\ell$ in
                $[-\ell/2,\ell/2)$. Show that $RE$ is, up to a global phase, a logical
                operator $\bar X^{m'}\bar Z^{m}$, and express the integers $m,m'$ in terms
                of $u$ and $v$. For which $(u,v)$ does the correction restore the original
                code state? Find the largest $r$ such that every shift with $|u|,|v|<r$ is corrected, and say how $r$ depends on $d$
            \end{enumerate}

  \item \emph{Noise level.} Let $u$ and $v$ be independent centred Gaussian variables of
        variance $\sigma^2$, as in the additive noise of \qref{Q023} with $\bar n=\sigma^2$.
        Write the probability $P_X$ that the correction of $u$ leaves a logical error as a
        series of Gaussian integrals, and find its leading behaviour as $\sigma\to0$. 
        
  \item \emph{Clifford gates are Gaussian.} Let $F$ be the Gaussian unitary with
        $F^\dagger\hat qF=-\hat p$ and $F^\dagger\hat pF=\hat q$ (a quarter turn of phase
        space), and on two modes let $C:=e^{-i\hat q_1\hat p_2}$.
        \begin{enumerate}[label=(\roman*)]
          \item Show that conjugation by $F$ maps $\bar X\mapsto\bar Z^{-1}$ and
                $\bar Z\mapsto\bar X$, as the qudit Fourier transform does, and that $C$
                acts on codewords as the logical sum gate
                $\ket{\bar\jmath_1,\bar\jmath_2}\mapsto\ket{\bar\jmath_1,\overline{j_1+j_2}}$.
          \item Show that for even $d$ the Gaussian unitary $e^{i\hat q^{2}/2}$ acts on
                the code as the logical phase gate $\ket{\bar\jmath}\mapsto
                e^{i\pi j^{2}/d}\ket{\bar\jmath}$. Show that for odd $d$ it does not
                preserve the code, and find the linear term $c\,\hat q$ that must be added
                to the exponent to obtain a logical diagonal gate.
        \end{enumerate}

  
\end{enumerate}
\end{problem}

\begin{solution}[Part (a)]
Multiplying the argument of $Z(a)$ or $X(a)$ by $m$, is equivalent to raising it to the power of $m$, $Z(ma)=Z(a)^m$. Then $\bar Z^{d}=Z(d\ell)=S_p$. Analogously $\bar X^d=S_q$. 

Using the central commutator \eqref{q025:weyl}
\begin{itemize}
    \item  $S_pS_q = e^{id^2\ell^2} X(d\ell)Z(d\ell) = e^{i2\pi d} X(d\ell)Z(d\ell)= S_qS_p$
\item $
    S_p\bar X = e^{id\ell^2} X(\ell)Z(d\ell) = e^{i2\pi } X(\ell)Z(d\ell)= \bar X  S_p$
\item $ S_q\bar Z = e^{-id\ell^2} Z(\ell)X(d\ell) = e^{-i2\pi } Z(\ell)X(d\ell)= \bar Z  S_q $
\item $ \bar Z\bar X = e^{i\ell^2} X(\ell)Z(\ell) = \omega  \bar X \bar Z$
\end{itemize}
    
To see that $S_q,S_p$ stabilizes these states, and that $\bar Z, \bar X$ acts as the corresponding qudit generalized Pauli gates (noticing that $e^{i\ell^2}=\omega$)

\begin{itemize}
    \item  $S_p\ket{\bar j} =  \sum_{k\in\ZZ} Z(d\ell)\ket{(dk+j)\ell} = \sum_{k\in\ZZ} \omega^{d(dk+j)}\ket{(dk+j)\ell} = \ket{\bar {j}} $, since $\omega^d=1$
\item $S_q\ket{\bar j} =  \sum_{k\in\ZZ} X(d\ell)\ket{(dk+j)\ell} = \sum_{k\in\ZZ} \ket{(d(k+1)+j)\ell} = \ket{\bar {j}} $, since $\sum_{k\in\ZZ}f(k)=\sum_{k\in\ZZ}f(k+1)$.
\item $\bar Z\ket{\bar j} =  \sum_{k\in\ZZ} Z(\ell)\ket{(dk+j)\ell} = \sum_{k\in\ZZ} \omega^{dk+j}\ket{(dk+j)\ell} = \omega^{j}\ket{\bar {j}} $, since $\omega^d=1$
\item $\bar X\ket{\bar j} =  \sum_{k\in\ZZ} X(\ell)\ket{(dk+j)\ell} = \sum_{k\in\ZZ} \ket{(dk+j+1)\ell} = \ket{\overline{j+1}} $.
\end{itemize}

\end{solution}

\begin{solution}[Part (b)]
i) Using the commutation rule
\begin{equation}
\begin{split}
     S_pE= Z(d\ell)X(u)Z(v)= e^{idu\ell}ES_p \\
     S_qE= X(d\ell)X(u)Z(v)= e^{-idv\ell}ES_q
\end{split}
\end{equation}
such that the eigenvalues of $S_p,S_q$ associated with the eigenvector $E\ket{\bar j}$ are, respectively, $e^{idu\ell},e^{-idv\ell}$. To see that this only reveals information about $u,v \mod \ell$, notice that if $u \equiv \bar u \mod \ell$, then $u = \ell \, m' +\bar u $ for some integer $m'$, such that 
$$e^{idu\ell}=e^{id (\ell \, m + \bar u )\ell}= \omega^{md} e^{id \bar u\ell}= e^{id \bar u\ell}.$$
Analogously for $v = m\ell + \bar v$.

ii) Then, using the previous notation,
\begin{equation}
\begin{split}
     RE &= X(-\bar u)Z(-\bar v) X(u)Z(v)\\
     &= e^{-iu\bar v} X(u-\bar u) Z(v-\bar v)\\
     &= e^{-iu\bar v} X(m'\ell) Z(m\ell) \\
     &= e^{-iu\bar v} X(m'\ell) Z(m\ell)\\
     &= e^{-iu\bar v} \bar X^{m'} \bar Z^m
\end{split}
\end{equation}

From that, we can see that the correction restores the original state, up to a phase, whenever $\bar X^{m'}=S_q$ and $\bar Z^{m}=S_p$, i.e., $m ,m'\equiv 0 \mod d.$ (Notice in the code space $S^k_q=S_q;\quad S^k_p=S_p$)

Now, this correction works for any $\bar u, \bar v = [-\ell/2 , \ell/2)$, which is to say 
$$kd\ell- \ell/2=m'\ell- \ell/2\leq u < m'\ell +\ell/2=kd\ell +\ell/2$$ and analogously for $v$, which means the set of correctable errors are 
\begin{equation}
    C = \cup_{k \in \ZZ} [d\ell k-\ell/2 , d\ell k+\ell/2)
    \label{correctable set}
\end{equation}

Then, whenever $|u|,|v|< \ell/2$ to code is corrected. This means that as the dimension of the logical qudits increases, the radius decreases by the relation $r\propto d^{-1/2}$.
\end{solution}

\begin{solution}[Part (c)]
The probability that the error in $X$ is correctable (using the correctable set \eqref{correctable set}
\begin{equation}
    P(u \in C) = \frac{1}{\sigma \sqrt{2\pi}}\sum_{_{k \in \ZZ}} \int_{d\ell k-\ell/2}^{d\ell k+\ell/2} dx e^{-x^2/2\sigma^2}.
\end{equation}

Clearly, the leading order integral happens for $k=0$. For $k\geq0$,the size of the interval is $\ell$ and in the integrand is smaller than $e^{-(d\ell k-\ell/2)^2/(2\sigma^2)}$. So the integrals become exponentially small as $\sigma \rightarrow 0$, and the factor $(d\ell k-\ell/2)^2$ makes the decay faster the greater it is. The same follows for $k<0$, with a factor $(d\ell k+\ell/2)^2$. Thus
\begin{equation}
    P(u \in C) \approx \frac{1}{\sigma \sqrt{2\pi}} \int_{-\ell/2}^{\ell/2} dx e^{-x^2/2\sigma^2}= \frac{2}{\sqrt{\pi}} \int_{0}^{\ell/(2\sqrt{2}\sigma)} e^{-t^2} = \text{erf
}\left(\frac{\ell}{2\sqrt{2}\sigma}\right)
\end{equation}

This determines the probability of a logical error to occur by the relation
\begin{equation}
    P_X = 1 -P(u \in C)
\end{equation}


\end{solution}

\begin{solution}[Part (d)]

i)
Given that $F$ is unitary, it follows that 
$$F^\dagger\hat q^n F= \underbrace{(F^\dagger \hat qF) (F^\dagger \hat qF) \cdots (F^\dagger \hat qF)}_{n\ \text{factors}}=(-\hat p)^n$$ as well as 
$$F^\dagger\hat p^n F= \underbrace{(F^\dagger \hat pF) (F^\dagger \hat pF) \cdots (F^\dagger \hat p F)}_{n\ \text{factors}}=\hat q^n.$$

Then, noticing how it acts in the series expansion of the exponential, it must also act accordingly in the exponentials as
\begin{equation}
\begin{split}
     F^\dagger X(a)F= F^\dagger e^{-ia\hat p}F= e^{-ia\hat q} = Z(-a)\\
     F^\dagger Z(b)F= F^\dagger e^{ib\hat q}F= e^{-ib\hat p} = X(b)
\end{split}
\end{equation}
being $F^\dagger \bar XF = \bar Z^{-1}$ and $F^\dagger \bar ZF = \bar X$ sub cases for $a=b=\ell$.

For the two-mode gate acting on codewords
\begin{equation}
    \begin{split}
        C\ket{\bar j_1, \bar j_2} &= \sum_{k\in\ZZ} e^{-i\hat q_1 \hat p_2}\ket{(dk+j_1)\ell} \otimes \ket{\bar j_2} \\
        &= \sum_{k\in\ZZ} e^{-i(dk+j_1)\ell\hat  p_2}\ket{(dk+j_1)\ell} \otimes \ket{\bar j_2}\\
        &= \sum_{k,k'\in\ZZ} \ket{(dk+j_1)\ell} \otimes e^{-i(dk+j_1)\ell\hat  p_2} \ket{(dk'+j_2)\ell}\\
        &= \sum_{k,k'\in\ZZ} \ket{(dk+j_1)\ell} \otimes \ket{(d(k+k')+j_1+j_2)\ell}\\ 
        &= \sum_{k,k'\in\ZZ} \ket{(dk+j_1)\ell} \otimes \ket{(dk'+j_1+j_2)\ell}\\ 
        &= \ket{\bar j_1, \overline{j_1 +j_2}}
    \end{split}
\end{equation}

ii) 
\begin{equation}
    \begin{split}
        e^{i\hat q^2/2} \ket{\bar j} &= \sum_{k\in\ZZ} e^{i\hat q^2/2}\ket{(dk+j)\ell}\\
        &=  \sum_{k\in\ZZ} e^{i(dk+j)^2\ell^2/2}\ket{(dk+j)\ell}\\
        &=  \sum_{k\in\ZZ} \omega^{(dk+j)^2/2}\ket{(dk+j)\ell}\\
        &=  \omega^{j^2/2} \sum_{k\in\ZZ} \omega^{d^2k^2/2} \omega^{dk j}\ket{(dk+j)\ell}
    \end{split}
\end{equation}

Notice that $\omega^{dk j}=e^{2\pi k j}=1$ for any $j,k \in \ZZ$. However, 
\begin{equation}
    \omega^{d^2k^2/2}=e^{i\pi d k^2} =\begin{cases}
        1, & d\ \text{is even}\\
        (-1)^k,& d\ \text{is odd} 
    \end{cases}
\end{equation}
such that we get the desired expression for $d$ even, and that is does preserve the code for $d$ odd, due to the introduction of the alternating sign. So the linear term that must be added is precisely the one leading to an alternating sign.

Applying the linear term on the corresponding eigenstate of $\hat q$
\begin{equation}
    e^{ic\hat q} \ket{(dk+j)\ell} =  e^{icj\ell} e^{icdk\ell}  \ket{(dk+j)\ell}
\end{equation}
it follows that $e^{icj\ell}$ is just a constant phase, and the only term that matters is $e^{icdk\ell}$. We want it to be equal to $(-1)^k=e^{i\pi k}$, so possible choice is $cd\ell=\pi$, to cancel the alternating phase, resulting in 
\begin{equation}
    e^{i\hat q^2/2} e^{i\pi \hat q /d\ell} \ket{\bar j} = \omega^{j^2/2} \omega^{j/2} \ket{\bar j} 
\end{equation}
for $d$ odd.

\end{solution}



% \begin{remark}
% \end{remark}

\begin{takeaway}
\end{takeaway}

%  Everything below is optional.  An entry with no attempt block is complete.
%  Record only what would change what you do next time: typos and arithmetic
%  slips you caught yourself are artefacts of the medium, not facts about the
%  mathematics.  Add "public," to show this block in the public build.
% \attempt{
%   date       = ,               % ISO, YYYY-MM-DD -- required if the block is used
%   time       = ,
%   difficulty = ,               % 1-5, perceived
%   howitwent  = ,
%   insight    = ,
%   note       = ,
% }

\end{entry}